% ── Section 1: Introduction ────────────────────────────────────────────────────
\section{Introduction}
\label{sec:intro}

Bangladesh has a population of approximately 170~million people, the vast
majority of whom speak Bengali (Bangla) as their first language.
However, standard written Bengali---as used in formal media and education---
differs substantially from the \emph{regional spoken dialects} used in everyday
conversation.
At least 15 regionally distinct dialect groups have been identified across
Bangladesh, each with unique phonological, lexical, and prosodic characteristics
that diverge sharply from the standard Dhaka dialect used in broadcast
media~\citep{chowdhury2014bengali}.

A second, equally important phenomenon is \emph{code-switching} (CS): the
practice of mixing two languages within a single utterance.
Among educated, urban Bangladeshi speakers---particularly those aged 18--35---
intra-sentential Bengali--English CS is the norm rather than the exception.
A typical utterance from a Chittagong-dialect YouTube creator might be:
\begin{quote}
\small\textit{``আমি এই phone-টা unboxing করতাছি, honestly speaking এইটার
camera performance অনেক ভালো।''}\\
\normalsize (\emph{I am unboxing this phone, honestly speaking its camera
performance is very good.})
\end{quote}
Here, Bengali dialect words (\textit{করতাছি}) co-occur with standard Bengali
(\textit{অনেক ভালো}) and English technical terms (\textit{unboxing, camera,
performance}), all within one sentence.

\paragraph{The gap in existing resources.}
Existing Bengali ASR resources address either dialect \emph{or} code-switching,
but never both simultaneously.
\citet{ben10_2025} released Ben-10, a 78-hour corpus covering 10 Bengali
dialects---but with no English code-switching, achieving a best WER of 70\%
on unseen dialects.
\citet{mucs2021} released the MUCS~2021 challenge dataset with Bengali--English
CS speech---but only in standard Dhaka Bengali, with no dialectal variation.
The intersection---\emph{dialectal Bengali mixed with English}---remains
entirely uncovered (Table~\ref{tab:gap}).

\begin{table}[h]
\centering
\caption{Existing Bengali ASR datasets and the gap filled by BanglaMix.}
\label{tab:gap}
\small
\begin{tabular}{lccc}
\toprule
\textbf{Dataset} & \textbf{Dialect?} & \textbf{CS?} & \textbf{Hours} \\
\midrule
OpenSLR-37 \citep{kjartansson2018bengali}   & \ding{55} & \ding{55} & 100 \\
MUCS 2021 \citep{mucs2021}                  & \ding{55} & \ding{51} & 16  \\
Ben-10 \citep{ben10_2025}                   & \ding{51} & \ding{55} & 78  \\
\midrule
\textbf{BanglaMix (ours)}                   & \ding{51} & \ding{51} & \textbf{164} \\
\bottomrule
\end{tabular}
\end{table}

\paragraph{Contributions.}
This paper makes the following contributions:
\begin{enumerate}[leftmargin=*, itemsep=2pt]
  \item \textbf{BanglaMix dataset}: 51,857 clips / 164.74~hours of dialectal
        Bengali--English CS speech across 15 regional dialects, sourced from
        YouTube and Facebook using dialect-name keyword queries.
  \item \textbf{Automated pipeline}: an end-to-end data collection pipeline
        using VAD segmentation, Whisper-based silver-standard transcription,
        automatic CS detection via Unicode script analysis, and a single
        strong multilingual LLM judge (Qwen2.5~\citep{qwen2025}) for borderline
        transcript validation---enabling large-scale annotation without any
        human transcribers.
  \item \textbf{LLM-ensemble evaluation}: beyond surface WER, this work
        introduces a three-model LLM-as-a-judge ensemble (Mistral-7B,
        LLaMA-3.1-8B, and Gemma) that rates ASR hypotheses for semantic
        correctness, code-switch preservation, and dialectal naturalness,
        capturing errors that string-level metrics miss.
  \item \textbf{Comprehensive benchmark}: systematic evaluation of four ASR
        architectures (Whisper, MMS, wav2vec2-XLS-R, WavLM) under
        zero-shot and fine-tuned conditions, with metrics including WER, CER,
        MIX-ER, and CS-detection F1.
  \item \textbf{Ablation study}: we show that training on the combination of
        dialect-only and code-switched clips outperforms either subset alone,
        and that LM shallow fusion further improves CTC model performance.
  \item \textbf{Cross-dialect generalisation}: per-dialect WER analysis
        demonstrating which dialect--model combinations pose the greatest
        remaining challenge.
\end{enumerate}

\paragraph{Paper structure.}
Section~\ref{sec:related} reviews related work.
Section~\ref{sec:dataset} describes dataset construction.
Section~\ref{sec:method} describes the ASR models and training procedure.
Section~\ref{sec:experiments} presents the experimental setup.
Section~\ref{sec:results} reports results.
Section~\ref{sec:analysis} provides error analysis and cross-dialect
generalisation results.
Section~\ref{sec:conclusion} concludes.
